\section{Related Work}
\label{sec:related}
\subsection{LEO Network Measurements}
Early measurements of Starlink reported median RTTs of a few tens of milliseconds with frequent spikes and loss under load~\cite{michel2022first}. Later studies traced many latency jumps to a global scheduler that reassigns users to satellites every 15\,s~\cite{mohan2024multifaceted,tanveer2023making}. Long campaigns characterized the access, point-of-presence, and backbone segments~\cite{pan2023measuring}, and Starlink throughput was shown to vary on time scales from sub-second to diurnal~\cite{garcia2023multitimescale}. OneWeb has no inter-satellite links in service, and its latency changes at every satellite and gateway handover~\cite{zhao2025measuring}. A study on the same testbed as ours measured OneWeb dynamics and the benefit of combining paths~\cite{perrin2025measurement}. These works explain \emph{why} LEO link state changes quickly. They do not ask which part of that change can be forecast at the sub-second scale of a packet scheduler, which is the question we answer.

\subsection{Network Performance Prediction}
Learned prediction of throughput is well established for cellular links. Lumos5G maps location, mobility, and radio context to 5G throughput~\cite{narayanan2020lumos5g}. A study of sub-second cellular throughput found that learned models gain little over simple predictors at fine time scales~\cite{basit2024predictability}. For LEO networks, recent work predicts Starlink throughput from months of measurements~\cite{liu2025vivisecting} and global Starlink latency and throughput at coarse resolution~\cite{benghe2026horizon}. LEO RTT has been predicted with spatiotemporal attention~\cite{tian2025attention}, and a concurrent preprint forecasts LEO link state with physics-informed features and quantile bands~\cite{chi2026reading}. These works evaluate point accuracy, and where intervals are reported they are not checked for coverage on sessions held out from calibration. None of them examines whether the input alignment of the measurement pipeline is causal. Simple linear forecasters are often competitive with deep sequence models~\cite{zeng2023transformers}, and we therefore compare against linear and naive references throughout.

\subsection{Distributional Forecasting and Calibration}
A proper scoring rule rewards a forecaster for reporting its true belief~\cite{gneiting2007strictly}, and the best point forecast depends on the loss used to judge it~\cite{gneiting2011making}. Quantile regression estimates conditional quantiles with the pinball loss~\cite{koenker1978regression}, and mean-variance networks estimate a Gaussian mean and variance~\cite{nix1994estimating}. Deep ensembles improve both accuracy and uncertainty~\cite{lakshminarayanan2017simple}. Split conformal prediction gives finite-sample coverage when calibration and test data are exchangeable~\cite{lei2018distribution}, and conformalized quantile regression (CQR) makes the interval width depend on the input~\cite{romano2019conformalized}. Conformal methods have been applied to multi-horizon recurrent forecasts~\cite{stankeviciute2021conformal} and extended to non-exchangeable data with explicit coverage gaps~\cite{barber2023conformal}. Adaptive versions track distribution shift online~\cite{gibbs2021adaptive}.

\subsection{Volatility and Burstiness}
In econometrics, a series whose level is hard to forecast but whose variance is predictable is modeled by autoregressive conditional heteroscedasticity~\cite{engle1982autoregressive}. Network traffic has long been known to be bursty at many time scales~\cite{paxson1995wide}. We find a related picture for sub-second LEO and 5G RTT. Its level is hard to forecast, while its spread depends on a link state that a few seconds of input reveal, and we use this to explain where learning helps.

\subsection{Multipath Scheduling}
Multipath QUIC (MPQUIC) extends QUIC to several paths~\cite{deconinck2017multipath}. Its default scheduler sends on the path with the lowest smoothed RTT, a policy inherited from MPTCP~\cite{paasch2014experimental}. BLEST and ECF account for head-of-line blocking and completion time on heterogeneous paths~\cite{ferlin2016blest,lim2017ecf}. MAMS uses mobility and handover awareness~\cite{yang2024mams}, and path-state forecasts have been fed to a 5G multi-access scheduler~\cite{maglione2025intelligent}. Measured MPTCP and MPQUIC sessions over Starlink and cellular paths show that these schedulers struggle with LEO dynamics~\cite{baltaci2023multipath}. All of them act on point estimates of path state.

\subsection{Positioning}
Three elements distinguish this paper. First, an analysis of sub-second measurements of two LEO systems and a 5G link that separates the predictable part of link state (its spread) from the hard part (its level). Second, a causality audit of the measurement pipeline with a model-level estimate of the resulting bias. Third, a forecaster whose uncertainty is calibrated on held-out sessions and evaluated for \emph{conditional} coverage and for its value to a deadline-aware multipath scheduler, against Gaussian, linear, and volatility-scaled alternatives.
